\documentclass[10pt,a4paper]{amsart}
\usepackage{amsmath,amssymb,amsthm,graphicx}
% For commutative diagrams we use Xy-pic
\usepackage[all]{xy}

% We use 2cell for 2-commutative diagrams.
\xyoption{2cell}
% Theorem environments.
%
\theoremstyle{plain}
\newtheorem{theorem}[subsection]{Theorem}
\newtheorem{proposition}[subsection]{Proposition}
\newtheorem{lemma}[subsection]{Lemma}

\theoremstyle{definition}
\newtheorem{definition}[subsection]{Definition}
\newtheorem{example}[subsection]{Example}
\newtheorem{exercise}[subsection]{Exercise}
\newtheorem{situation}[subsection]{Situation}

\theoremstyle{remark}
\newtheorem{remark}[subsection]{Remark}
\newtheorem{remarks}[subsection]{Remarks}

\numberwithin{equation}{subsection}

% Macros
%
\def\lim{\mathop{\mathrm{lim}}\nolimits}
\def\colim{\mathop{\mathrm{colim}}\nolimits}
\def\Spec{\mathop{\mathrm{Spec}}}
\def\Hom{\mathop{\mathrm{Hom}}\nolimits}
\def\Ext{\mathop{\mathrm{Ext}}\nolimits}
\def\SheafHom{\mathop{\mathcal{H}\!\mathit{om}}\nolimits}
\def\SheafExt{\mathop{\mathcal{E}\!\mathit{xt}}\nolimits}
\def\Sch{\mathit{Sch}}
\def\Mor{\mathop{\mathrm{Mor}}\nolimits}
\def\Ob{\mathop{\mathrm{Ob}}\nolimits}
\def\Sh{\mathop{\mathit{Sh}}\nolimits}
\def\NL{\mathop{N\!L}\nolimits}
\def\CH{\mathop{\mathrm{CH}}\nolimits}
\def\proetale{{pro\text{-}\acute{e}tale}}
\def\etale{{\acute{e}tale}}
\def\QCoh{\mathit{QCoh}}
\def\Ker{\mathop{\mathrm{Ker}}}
\def\Im{\mathop{\mathrm{Im}}}
\def\Coker{\mathop{\mathrm{Coker}}}
\def\Coim{\mathop{\mathrm{Coim}}}

% Boxtimes
%
\DeclareMathSymbol{\boxtimes}{\mathbin}{AMSa}{"02}

%
% Macros for moduli stacks/spaces
%
\def\QCohstack{\mathcal{QC}\!\mathit{oh}}
\def\Cohstack{\mathcal{C}\!\mathit{oh}}
\def\Spacesstack{\mathcal{S}\!\mathit{paces}}
\def\Quotfunctor{\mathrm{Quot}}
\def\Hilbfunctor{\mathrm{Hilb}}
\def\Curvesstack{\mathcal{C}\!\mathit{urves}}
\def\Polarizedstack{\mathcal{P}\!\mathit{olarized}}
\def\Complexesstack{\mathcal{C}\!\mathit{omplexes}}
% \Pic is the operator that assigns to X its picard group, usage \Pic(X)
% \Picardstack_{X/B} denotes the Picard stack of X over B
% \Picardfunctor_{X/B} denotes the Picard functor of X over B
\def\Pic{\mathop{\mathrm{Pic}}\nolimits}
\def\Picardstack{\mathcal{P}\!\mathit{ic}}
\def\Picardfunctor{\mathrm{Pic}}
\def\Deformationcategory{\mathcal{D}\!\mathit{ef}}
\title{An intermediate constant subsheaf of a constant sheaf on the etale site of an irreducible scheme is constant, in the set, abelian group and module cases}
\author{The Stacks Project Authors}
\date{stacks-project:etale-cohomology.tex:local-snapshot-20261011; GNU FDL 1.2 or later}
\begin{document}
\maketitle
\noindent\textit{Editorial note.} All delivered mathematical text is the author's own, copied from the pinned Stacks Project chapter file etale-cohomology.tex; no statement or proof step was rewritten, reordered or removed.\par\medskip
\noindent\textit{Editorial note.} The counted proof is the author's complete argument: The counted proof is the author's complete argument: the one-sentence reason is the density of the generic point in every nonempty etale neighbourhood, which is exactly the content of the claim.\par\medskip
\noindent\textit{Editorial note.} The label lemma-irreducible-subsheaf-constant-zero defines the statement; the counted proof is the proof environment immediately following the lemma in the source.\par\medskip
\noindent\textit{Editorial note.} The preamble is the chapter's own native preamble excerpt already used for this source: the Xy-pic setup (native preamble lines 14-18) and the native theorem declarations with the mathematical macros (native preamble lines 154-217). The class block and the externaldocument lines of the native preamble are not included, because the assembler supplies the wrapper class and those external documents are not part of this material.\par\medskip
\noindent\textit{Allowed standard prerequisites (editorial).}\begin{itemize}
\item The constant sheaf underline M on X\_etale and its stalks. Same chapter, the constant sheaf underlying a set / abelian group / module and the geometric-point stalks used implicitly in the conclusion G = underline\{S'\} . Supplies the objects S', S / A', A / M', M between which the sheaf G is sandwiched and the value of the stalk over the generic point eta.
\item The generic point eta lies in the image of every etale morphism. Same chapter, the density of the generic point of the irreducible scheme X under etale morphisms with nonempty source. Supplies the single fact used in the proof: for every etale U -> X with U nonempty, eta is in the image.
\end{itemize}\medskip
\begin{lemma}
\label{lemma-irreducible-subsheaf-constant-zero}
Let $X$ be an irreducible scheme with generic point $\eta$.
\begin{enumerate}
\item Let $S' \subset S$ be an inclusion of sets. If we have
$\underline{S'} \subset \mathcal{G} \subset \underline{S}$
in $\Sh(X_\etale)$ and $S' = \mathcal{G}_{\overline{\eta}}$, then
$\mathcal{G} = \underline{S'}$.
\item Let $A' \subset A$ be an inclusion of abelian groups. If we have
$\underline{A'} \subset \mathcal{G} \subset \underline{A}$
in $\textit{Ab}(X_\etale)$ and $A' = \mathcal{G}_{\overline{\eta}}$, then
$\mathcal{G} = \underline{A'}$.
\item Let $M' \subset M$ be an inclusion of modules over a ring $\Lambda$.
If we have $\underline{M'} \subset \mathcal{G} \subset \underline{M}$
in $\textit{Mod}(X_\etale, \underline{\Lambda})$
and $M' = \mathcal{G}_{\overline{\eta}}$, then
$\mathcal{G} = \underline{M'}$.
\end{enumerate}
\end{lemma}

\begin{proof}
This is true because for every \'etale morphism $U \to X$
with $U \not = \emptyset$ the point $\eta$ is in the image.
\end{proof}

\end{document}
